\chapter{Binge Eating}

\section*{Definition}
Binge eating is the consumption of an objectively large amount of food within a discrete period (e.g., 2 hours) accompanied by a sense of loss of control. For binge-eating disorder (BED), episodes occur at least once weekly for 3 months and are associated with marked distress.

\section*{When to See a Doctor}

\begin{itemize}
 \item Recurrent binge episodes with loss of control and distress, even if they do not yet meet the formal frequency or duration criteria for binge-eating disorder
 \item Eating is accompanied by extreme distress, shame, or impaired functioning
 \item There is self-induced vomiting, misuse of laxatives or diuretics, fasting, excessive exercise, significant weight change, or medical complications
 \item There are thoughts of self-harm, severe depression, fainting, chest pain, or other urgent physical or mental-health concerns
\end{itemize}

\section*{How It Develops}
Binge eating usually develops through interacting biological, psychological, and social factors. Hunger from irregular or restrictive eating, emotional distress, learned responses to food cues, and difficulty regulating emotions can contribute. The mechanisms are complex and are not explained by a single brain pathway or lack of willpower.

\section*{Causes}
\begin{itemize}
 \item \textbf{Psychiatric:} Binge-eating disorder, bulimia nervosa, or another eating disorder
 \item \textbf{Psychological factors:} Negative affect, emotional dysregulation, impulsivity, low self-esteem, body dissatisfaction, dietary restraint
 \item \textbf{Biological and genetic:} Family history and individual differences in appetite, reward, and stress regulation
 \item \textbf{Environmental:} Childhood obesity, history of weight-based teasing, trauma, food insecurity
 \item \textbf{Medical or medication-related:} Conditions or medicines that alter appetite, sleep, mood, or impulse control; rare neurological disorders are uncommon causes
\end{itemize}

\section*{Related Symptoms}
\begin{itemize}
 \item Obesity, emotional eating, food cravings, body dissatisfaction, depression, anxiety, substance use disorders
\end{itemize}
\textbf{Important:} Binge-eating disorder can occur at any body size and commonly co-occurs with depression, anxiety, trauma-related symptoms, and metabolic disease. Body weight alone does not diagnose it.

\section*{Medical Evaluation}
\begin{itemize}
 \item Clinical interview assessing frequency of binge episodes, loss of control, distress, and compensatory behaviors
 \item Screen for comorbid psychiatric conditions: depression, anxiety, substance use
 \item Physical assessment considers pulse, blood pressure, weight trends, hydration, nutritional status, and complications without using weight as the sole measure of severity.
 \item Blood tests such as glucose, lipids, or liver tests are selected according to symptoms, medicines, risk factors, and physical findings.
 \item Screening questionnaires can support assessment, but diagnosis requires a clinical interview; ask specifically about vomiting, laxatives, fasting, excessive exercise, mood, substance use, and self-harm.
\end{itemize}

\section*{Treatment Options}
\begin{itemize}
 \item First-line psychological treatment is an eating-disorder-focused guided self-help program; if it is unsuitable or ineffective, CBT focused on binge eating or interpersonal therapy may be offered.
 \item Lisdexamfetamine is approved for moderate-to-severe BED in adults in the United States, but it is not suitable for everyone and requires medical screening and monitoring. Antidepressants may be considered for BED or coexisting depression/anxiety; medicines are not a substitute for psychological treatment.
 \item Nutritional support focuses on regular, adequate meals and reducing restrictive eating. Weight loss is not the immediate treatment target, and dieting can trigger binge episodes.
 \item Metabolic conditions and weight-related health concerns should be addressed sensitively and separately, without compromising eating-disorder treatment.
\end{itemize}

\section*{Self-Care and Home Management}
\begin{itemize}
 \item Keep a food-mood diary to identify triggers and patterns
 \item Practice mindful eating: eat slowly, without distractions, and pause between bites
 \item Establish regular meal patterns (3 meals + 1--2 planned snacks) to reduce deprivation
 \item Avoid rigid food rules and dieting that promote restriction-binge cycles
\end{itemize}

\section*{References}
\begin{thebibliography}{99}
\bibitem{binge_eating:ref1} Brownley KA, Berkman ND, et al. ``Binge-eating disorder in adults: a systematic review and meta-analysis.'' \textit{Ann Intern Med}. 2016;165(6):409-420.
\bibitem{binge_eating:ref2} McElroy SL, Hudson JI, et al. ``Pharmacologic treatments for binge-eating disorder.'' \textit{J Clin Psychiatry}. 2020;81(3):OT19048COM.
\bibitem{binge_eating:ref3} Treasure J, Duarte TA, et al. ``Eating disorders.'' \textit{Lancet}. 2020;395(10227):899-911.
\bibitem{binge_eating:ref4} Sadock BJ, Sadock VA, et al. \textit{Kaplan \& Sadock\'s Comprehensive Textbook of Psychiatry}, 10th ed. Wolters Kluwer, 2017.
\bibitem{binge_eating:ref5} American Psychiatric Association. ``The American Psychiatric Association Practice Guideline for the Treatment of Patients With Eating Disorders.'' 2023. https://www.psychiatry.org/getmedia/97405f0d-1bd4-43d0-abdd-c013fcd8686d/APA-Eating-Disorders-Practice-Guideline-Under-Copyediting.pdf.
\bibitem{binge_eating:ref6} National Institute for Health and Care Excellence. ``Eating disorders: recognition and treatment (NG69).'' Updated 2025. https://www.nice.org.uk/guidance/ng69.
\end{thebibliography}
